\begin{textsample}{POS Dim 1 – ce – Score 51.00 – ce/ce\_em\_52.txt}  \label{ex:f1_pos_048}
8.4.8 Saberes, fazeres, práticas e vivências através de exercícios e avaliações internas e externas para produção de conhecimento significativo A partir do início dos anos 1980, com a redemocratização do ensino e com o movimento pela profissionalização da atividade docente, foi reconhecida a existência e necessidade de saberes e práticas específicas a profissão do professor.
Com isto, houve \textbf{um} crescimento de pesquisas sobre a temática, tendo a formação de professores e estes saberes específicos como objeto de estudo, especialmente na década de noventa.
Os documentos oficiais da Unesco \textbf{que} tratam da questão educacional, dentre eles, o Relatório da Comissão Internacional sobre Educação para o \textbf{século} XXI ( 1996 ), lançaram luz sobre a questão e proporcionou \textbf{um} debate mundial sobre o campo educacional e as necessidade de saberes aplicados a educação, a partir da escrita dos quatro pilares da educação ou quatro aprendizagens fundamentais, considerados também os pilares do conhecimento : aprender a conhecer, aprender a fazer, aprender a \textbf{viver} junto, aprender a ser.
Através destes, o indivíduo poderá, respectivamente, “ adquirir os instrumentos da compreensão ”, “ poder agir sobre o meio envolvente ”, “ participar e cooperar com os outros em todas as atividades humanas ” e através do quarto pilar ter a representação da “ via essencial \textbf{que} integra as três precedentes ”.
Levando -se em conta ainda a multiplicidade de pontos de encontro, de relação e de trocas entre estes elementos, \textbf{que} constituem \textbf{uma} única via, a educacional., fortalece o debate sobre esses aspectos, quando descreve os sete saberes necessários à educação do futuro, nos quais aborda questões sobre vários níveis de ensino, apresentando os “ buracos negros da educação ”, até então ignorados, ou fragmentados dentro dos programas educativos \textbf{que} atendem á Educação Básica e o Ensino Superior.
Nesta proposta, o primeiro ponto discutido é o conhecimento, pois este é o produto essencial do ensino.
Ele dá luz à necessidade de saber a importância de conhecer, abordando também a ilusão, ou como somos iludidos sobre o mundo e sobre nossa realidade dentro deste “ saber ”.
O \textbf{autor}, a partir dessa premissa, apresenta e discute, as lacunas e saberes \textbf{que} devem condicionar a educação e o ensino do futuro. a ) \textbf{Um} conhecimento capaz de criticar o próprio conhecimento ; b ) Discernir as informações chave, tendo claros os princípios do conhecimento pertinente ; c ) Ensinar a condição humana ; d ) Ensinar a identidade terrena ; e ) Enfrentar as incertezas ; f ) Ensinar a compreensão ; g ) A ética do gênero humano.
Esses saberes, segundo Morin ( 2000, p.
XX ), possibilitaria \textbf{que} criássemos caminhos para alcançar o objetivo da educação, entendendo \textbf{que} “ o dever principal da educação é preparar cada \textbf{um} para enfrentar os não saberes com lucidez", fazendo com \textbf{que} o indivíduo perceba a necessidade de “ integrar os \textbf{erros} nas concepções para \textbf{que} o conhecimento \textbf{consiga} avançar ”.
A integração de saberes, práticas e vivências nos \textbf{remete} à reflexão sobre a formação docente, sobre a formação continuada, enfim, toda a formação do professor ao longo da vida.
Discutir esta temática envolve quatro aspectos : formação técnico-científica, formação prática, formação pedagógica e formação política.
Estes \textbf{tópicos} contemplam : a busca por atualização constante, para evitar atitudes repetitivas na atividade docente ; a necessidade da relação entre \textbf{teoria} e prática ; desenvolvimento de competências e habilidades para a \textbf{docência} ; experiência profissional ; planejamento de ensino ; conhecimentos sobre aluno, instituição e seus objetivos ; avaliações, seleção de conteúdos, das atividades e dos recursos ; além de construções e reconstruções do conhecimento, etc. ( SOUSA, 2010 ).
Valcemiro Nossa ( 1999 ) enfatiza \textbf{que} a formação do docente vai muito além de conhecimentos técnicos \textbf{exigidos} e propriamente ditos, mas deve abranger \textbf{uma} gama de conhecimentos filosóficos, éticos, políticos, entre outros.
Outro especialista da área da educação, Tardif ( 2004, p. 54 ) afirma \textbf{que} o saber docente é \textbf{um} “ saber plural, formado de diversos saberes provenientes das instituições de formação, da formação profissional, dos currículos e da prática cotidiana ”, destaca a existência de quatro tipos de saberes para a atividade dos professores, conforme \textbf{vemos} a seguir : 1.
Saberes da Formação Profissional : são \textbf{baseados} nas ciências e na erudição, o professor adquire durante sua formação inicial e/ou continuada, os quais os conhecimentos ou saberes pedagógicos ou práticos, \textbf{métodos} e técnicas do saber fazer docente, “ legitimados cientificamente e igualmente transmitidos aos professores ao longo do seu processo de formação ”. 2.
Saberes Disciplinares : são os saberes produzidos socialmente e \textbf{historicamente} pela \textbf{humanidade}, administrados academicamente pelas instituições e comunidade científica, pertencentes a campos diversos do conhecimento.
Ex. : ciências exatas, ciências humanas, linguagem, etc. 3.
Saberes Curriculares : estes conhecimentos estão relacionados aos modelos e formas como as instituições educativas administram os saberes e como estes são trabalhados com seus alunas/os ( saberes disciplinares ). “ Apresentam -se, concretamente, sob a forma de programas escolares ( objetivos, conteúdos, \textbf{métodos} ) \textbf{que} os professores devem aprender e aplicar ”. 4.
Saberes Experienciais : São próprios do exercício da atividade profissional dos docentes.
São saberes produzidos pelos docentes por meio de suas vivências em situações específicas relacionadas ao ambiente escolar e às relações com alunas/os e colegas profissionais.
Nesse sentido, “ incorporam -se à experiência individual e coletiva sob a forma de habitus e de habilidades, de saber-fazer e de saber ser ” ( p. 38 ).
A integração de saberes é necessária e se fortalece através da formação de professores, a qual \textbf{revela} em seu caráter prático a importância da relação entre \textbf{teoria} e prática diante dos saberes científicos.
Montero ( 2001 ) relata \textbf{que} : > Quando a relação teoria-prática se transforma numa passarelle onde se exibem as últimas criações da moda, ou em slogans, no lugar dos princípios operativos, não só está a pôr em risco a sua identidade como campo disciplinar perigoso e concernente a \textbf{um} conjunto de saberes em forma de saber fazer corre o perigo de se converter num conjunto de despropósitos.
É importante \textbf{que} nas ciências humanas as Tecnologias da Informação sejam utilizadas para potencializar o aprendizado dos \textbf{discentes}, haja vista \textbf{que} grande parte dos objetos de conhecimento da área dialoga com essas tecnologias.
Assim, o desenvolvimento da prática docente na área das ciências humanas deve ter como \textbf{fundamento} a formação de professores em sentido integral, incluindo o conhecimento e uso das TIC na educação, como nos \textbf{diz} Sousa ( 2010 ) : > O uso das TIC deve ser entendido como estratégia de inovação, mantendo os \textbf{compromissos} com os fins da ação educativa, pois para responder às demandas quantitativas e qualitativas da educação faz -se indispensável buscar novas estratégias, \textbf{uma} vez \textbf{que} os modelos convencionais não têm \textbf{conseguido} responder aos desafios impostos pela expansão quantitativa e pela exigência qualitativa.
A formação docente, seus saberes, práticas e vivências formam \textbf{um} conjunto importante para \textbf{que} o professor possa executar seu trabalhar com qualidade, além de produzir instrumentos eficazes, tais como exercícios e avaliações para a produção de \textbf{um} conhecimento significativo aos alunas/os e para ele próprio, dentro do processo ensino aprendizagem.
Embora, em muitos países pelo mundo, já existem modelos de escolas \textbf{que} experimentam não utilizar a avaliação, com saldo positivo sobre o processo ensino aprendizagem, na educação brasileira, este instrumento se faz presente de maneira \textbf{forte} e torna -se necessário para \textbf{uma} série de processos pelo qual o \textbf{educando} deve passar, acompanhando -o do Ensino Infantil ao Ensino Superior.
Nesse sentido, pensar e debater sobre avaliação é essencial, pois perpassa o contexto do Ensino Médio e da própria vida educacional, com resultados \textbf{que} vão refletir socialmente e culturalmente na vida do indivíduo.
Quanto à avaliação, entende -se \textbf{que} deve ser processual e contínua percorrendo todo o caminho \textbf{que} vai desde o planejamento da aula até a execução.
A avaliação deve ser pensada além de \textbf{um} instrumento de medição cognitiva do aluno, contemplando também o modo operante do professor na relação ensino e aprendizagem, porquanto, a \textbf{atual} organização do trabalho pedagógico foi alterada com a inclusão de novas tecnologias e de novos espaços \textbf{que} surgem com ela.
A avaliação não pode constituir -se como \textbf{um} instrumento punitivo ou classificatório, \textbf{que} tem como principal consequência a diminuição da autoestima das/os alunas/os, reforçando a competitividade e o individualismo.
A avaliação como \textbf{um} instrumento de desenvolvimento cognitivo deve promover a/o aluna/o e não suscitar medo e angústia.
Priorizar o seu aprendizado \textbf{exige} de nós docentes estratégias positivas de avaliação \textbf{que} colaborem para o desenvolvimento de suas competências e habilidades.
As avaliações, portanto, devem ser desenvolvidas dentro de perspectivas \textbf{contemporâneas}, não tradicionais, buscando valorizar a compreensão dos acontecimentos humanos em sociedade a partir de \textbf{um} tempo histórico e \textbf{que} tem relação com o presente e o futuro.
Quanto ao ENEM, e demais avaliações externas, não se pode negar a importância destas, entretanto, não se deve permitir \textbf{que} a formação fornecida pela educação básica se torne sua refém, do contrário ela se desviará dos seus outros propósitos : formação para o mundo do trabalho e para o exercício da cidadania.
O exercício e avaliação por repetição, para decorar e com isso atender objetivos imediatistas como, por exemplo, vestibulares e concursos, não garante a construção de conhecimentos \textbf{sólidos} e não prepara alunas/os críticos, mas \textbf{meros} repetidores do discurso \textbf{hegemônico}, incapazes de pensarem de forma reflexiva e crítica.
Adotando o exercício de criticar os noticiários locais e nacionais \textbf{que} reproduzem notícias parciais sobre os cotidianos sociais, desprovidos de \textbf{fundamentos}, o professor proverá seu aluno de \textbf{um} importante elemento em \textbf{direção} a formação de cidadãos conscientes.
Enfim, pensar sobre o trabalho do professor da área de Ciências Humanas, assim como nas outras áreas, deve levar a reflexões sobre a necessidade de \textbf{um} \textbf{sólido} conhecimento \textbf{teórico} \textbf{que} fundamente a sua prática em sala de aula.
E esta prática docente está ligada às concepções de mundo do professor, como ele \textbf{vê} a si próprio, a sociedade na qual está inserido, sua busca pessoal dentro da profissão e os caminhos de investimentos nela.
A consciência do papel social do docente possibilita ao mesmo mapear, na escola, os potenciais físicos e humanos \textbf{que} poderão ser utilizados em sua prática \textbf{metodológica}.
As escolas cearenses estão inseridas em espaços e realidades sociais distintas e específicas, o \textbf{que} deve redobrar a atenção dos \textbf{que} promovem e \textbf{fomentam} o processo de ensino e aprendizagem.
O professor, por exemplo, deve sempre buscar conhecer o contexto em \textbf{que} está inserido para \textbf{que} em sua prática possa contemplar essas especificidades.
Escolas pertencentes à zona rural ou urbana, em locais de grandes conflitos sociais, com altos índices de marginalidade, violência, cujos problemas sociais vivenciados pela comunidade convergem para o interior da escola e requerem maior atenção.
O espaço geográfico, no qual a escola está localizada, é fator determinante para \textbf{que} a prática de ensino dos professores de humanas seja significativa na aprendizagem dos seus alunos/as.
O docente deve ter sensibilidade para saber quem é essa comunidade, sua história, seus problemas e a partir disso propor \textbf{um} ensino \textbf{que} suscite os interesses das alunas/os.
Esta \textbf{percepção} dará condições a \textbf{um} ensino contextualizado, através de metodologias \textbf{que} permitam a leitura e interpretação do passado a partir do presente, ensejando perspectivas para a construção de \textbf{um} futuro promissor, além de sinalizar na \textbf{direção} do uso dos recursos pedagógicos e didáticos \textbf{que} auxiliem na otimização do ensino.
Desse modo, os diversos espaços internos e externos à escola podem enriquecer \textbf{uma} proposta de ensino crítico e reflexivo.
Determinados ambientes, fora da sala de aula, podem proporcionar \textbf{uma} mudança significativa nas estratégias de ensino, especialmente a \textbf{que} tem como princípio fundamental a investigação a partir da pesquisa escolar, ou mesmo \textbf{uma} \textbf{abordagem} a partir de obras \textbf{literárias} ou o uso de dicionários e obras de arte.
Neste contexto, a biblioteca, a exemplo de outros espaços de aprendizagem, permite subsidiar a prática do professor em relação à aprendizagem das/os alunas/os.
Ao professor também cabe investigar quais referenciais bibliográficos se encontram disponíveis na biblioteca e \textbf{que} softwares podem ser acessados para o desenvolvimento da pesquisa em Humanas com os seus alunas/os.
Tais recursos servem para estimular no aluno a reflexão sobre o conhecimento já produzido, favorecendo o desenvolvimento da pesquisa escolar.
Ao mediar o conhecimento, o professor deve estar formando alunas/os investigadores.
Deve também preparar indivíduos capazes de problematizar a realidade a partir da apreensão e compreensão do conhecimento formal.
Revisitando a obra de Freire ( 1970 ), ao criar condições para o entendimento dessa realidade, é possível ir além e partir dessa compreensão para criar os meios de poder mudar, agir e transformar a sociedade.
O pensamento freiriano indica \textbf{que} o caminho \textbf{metodológico} não conduz apenas ao conhecimento e sim a \textbf{um} processo de consciência \textbf{que} se efetiva em mudança da realidade do \textbf{discente} e de sua comunidade, em \textbf{uma} ação constante e \textbf{dialética}.
Essa “ curiosidade cognitiva ”, exercitada por professores e alunas/os, deverá ser alimentada pela constante indignação frente ao contexto social gerando \textbf{posicionamentos}, críticas, enfim o exercício pleno do ato político, de dialogar e agir em prol das \textbf{vontades} coletivas.
O uso da biblioteca e dos materiais bibliográficos disponíveis oferecem possibilidades para \textbf{que} o professor contorne possíveis limitações pedagógicas no \textbf{que} tange ao uso do livro didático.
A compreensão e interpretação das ciências \textbf{que} compõem a área ocorrem com o desenvolvimento de alunas/os leitores e escritores, \textbf{uma} vez \textbf{que} estas habilidades constituem competências essenciais no contexto \textbf{atual} do ensino de Humanas.
O planejamento de ações junto à biblioteca deve funcionar como \textbf{uma} estratégia \textbf{que} fortaleça o ensino, a pesquisa e a \textbf{interdisciplinaridade}.
Só assim, proporcionará \textbf{um} ensino em \textbf{que} o conhecimento seja entendido a partir de \textbf{uma} leitura geral, superando o paradigma \textbf{atual} marcado pela \textbf{fragmentação}, favorecendo \textbf{um} ensino \textbf{que} não seja apenas o relato de fatos, mas \textbf{uma} construção produtiva, humanizada e política, comprometida com transformações na realidade social.
Pensar na otimização da biblioteca constitui -se numa \textbf{tarefa} conjunta entre professores e alunas/os, em sinergia com os demais recursos disponíveis como o laboratório de informática e a internet.
As possibilidades \textbf{que} são permitidas pelas tecnologias de comunicação e informação como mediadoras do ensino e da aprendizagem, também podem suscitar \textbf{uma} ampliação no espaço de aprendizagem para além da sala de aula.
No ensino mediado por essas tecnologias, cabe ao professor mapeá -las na escola, buscar compreender de \textbf{que} forma e quando elas serão utilizadas, além de \textbf{fomentar} os usos das tecnologias como ferramentas \textbf{que} proporcionam interações em rede e \textbf{que} ampliam o universo das informações.
O laboratório de informática deve ser concebido como \textbf{um} espaço de criação, pesquisa e construção do conhecimento, \textbf{que} juntamente com a internet, podem possibilitar a elaboração de material didático a ser utilizado tanto por \textbf{discentes} como professores.
O como fazer deve ser investigado por professores e alunas/os, de acordo com a realidade, tendo em vista \textbf{que} são \textbf{poderosas} ferramentas para a pesquisa e construção de conhecimento.
O uso da biblioteca e da sala de informática deve ser feito com planejamento, atendendo aos conteúdos programados, para \textbf{uma} ação pedagógica mais efetiva.
O ensino das ciências humanas e sociais aplicadas, pode, também, ser enriquecido com práticas \textbf{metodológicas} \textbf{que} utilizem recursos midiáticos ( áudio, vídeo e imagem ), acompanhados de visão interdisciplinar sobre estas diversas linguagens e \textbf{que} possibilitem à mediação e aquisição efetiva do conhecimento, proporcionando \textbf{uma} visão crítica à realidade social e conduzindo os \textbf{sujeitos} a se sentirem parte ativa desta sociedade.
A aula de campo constitui mais \textbf{um} recurso \textbf{metodológico} \textbf{que} possibilita a investigação da realidade local e sua inserção numa perspectiva mais larga, e vice-versa.
A partir de \textbf{uma} perspectiva mais global pode -se perceber a sociedade local a qual os alunas/os pertencem e desenvolver o senso crítico e reflexivo do aluno na elaboração e construção da identidade.
Para melhor visualizar o \textbf{que} se fala sobre estas formas e espaços, imagine \textbf{que} o professor de Humanas, em seu planejamento anual, ao trabalhar com a sociedade brasileira no \textbf{século} XIX, faça \textbf{uma} \textbf{abordagem} do cotidiano dessa realidade histórica, filosófica, geográfica e sociológica a partir de \textbf{um} enfoque \textbf{literário}.
Na elaboração do plano de aula, ele faz \textbf{um} levantamento na internet de sites \textbf{que} contemplem obras a serem utilizadas e organize \textbf{um} estudo crítico e analítico a partir da literatura e à luz de \textbf{uma} \textbf{abordagem} histórica.
Tendo em vista, desse modo, \textbf{que} há \textbf{uma} variedade de construções e narrativas \textbf{que} \textbf{permeiam} o componente curricular, mas é na diversidade de \textbf{abordagens} \textbf{que} se permite contemplar os múltiplos aspectos de \textbf{uma} necessária alteridade \textbf{que} compreenda a complexidade da aprendizagem nas ciências humanas e aquele \textbf{que} se estabelece com \textbf{equidade}, além de ética, inclusiva e significativa.
Quanto aos usos dos recursos pelo professor, \textbf{que} \textbf{fomentam} os rumos dados ao conhecimento produzido pelos alunas/os, através dos espaços e instrumentos de aprendizagem, \textbf{que} estão disponíveis na escola e fora dela, bem como nas diversas linguagens comunicacionais abordadas, se permitam aos alunas/os entenderem e firmarem o seu lugar no conjunto da sociedade.
O exercício da curiosidade como metodologia para o ensino das Humanas, pode ser \textbf{um} dos caminhos para \textbf{que} sejam superadas práticas anacrônicas com novas roupagens.
Na elaboração de pesquisas ou qualquer trabalho no qual sua feitura ultrapasse o ambiente escolar, é \textbf{preciso} monitoramento e regras previamente estabelecidas.
Não devem ser aceitas reproduções ou transcrições de livros e revistas especializadas.
Além disso, os recursos da informática, o “ copiar/colar ” de pesquisas na Internet, pode transformar qualquer trabalho em \textbf{mera} cópia, onde a originalidade se resumirá na capa com o nome do \textbf{autor}.
A contestação deste modelo \textbf{hegemônico} do ensino de Humanas constitui -se num exercício constante.
A partir da pesquisa escolar, pode -se ampliar a visão analítica dos aprendentes, reforçando a salutar ideia de \textbf{que} a transferência de conhecimento não deve ser meramente transmissiva.
O conhecimento deve ser, portanto, significativo para os alunas/os, \textbf{enquanto} sujeitos dessa sociedade. \textbf{Interdisciplinaridade}, \textbf{contextualização} e diversidade são posturas de \textbf{uma} prática \textbf{que} não pode ser pensada de forma isolada, mas sim como processo.
O professor não se torna interdisciplinar só porque comentou a Lei de Newton em sua aula, por exemplo.
Entretanto, se ele estimular a investigação, juntamente com seus alunas/os, sobre as implicações desse pensamento para a fundação de \textbf{um} novo conhecimento e de \textbf{uma} nova ordem social, com o apoio do professor de Física, então ele vislumbra, de maneira efetiva, a construção do conhecimento interdisciplinar em sua tentativa de explicar a sociedade.
Da mesma forma, é possível pensar a questão da diversidade étnica e racial.
Não se pode imaginar \textbf{que} ao falar do negro e do indígena, do europeu, do nordestino, se está propondo \textbf{um} ensino para a diversidade.
Os exercícios e as avaliações devem seguir \textbf{um} caminho \textbf{que} ajude a desconstruir os estereótipos e as visões distorcidas contidas no imaginário educacional, onde o negro é \textbf{visto} pela sua dança e suas comidas, o indígena como o nativo das Américas, o preguiçoso \textbf{que} não se adaptou ao trabalho agrícola no Brasil Colonial.
Agora ambos – negro e indígena – lutam por terra e direitos sociais, para \textbf{que}?
E o nordestino, cabeça chata, comedor de farinha \textbf{que} vai para o sul em busca de comida e de sobrevivência?
Desconstruir esse discurso se coloca como \textbf{um} exercício para \textbf{um} ensino na área de Ciências Humanas \textbf{que} visa conduzir a formação de indivíduos críticos, ativos e investigativos.
A curiosidade epistemológica deve provocar \textbf{uma} revolução na prática do professor para \textbf{que} ele pense na pesquisa como \textbf{uma} grande aliada pedagógica, principalmente, se ele articular, de modo eficaz, o uso do computador, da internet, da biblioteca, dos museus, das ruas, dos centros históricos, da memória da escola, da cidade, dos \textbf{sujeitos} envolvidos da comunidade, enfim, \textbf{um} conjunto de elementos \textbf{que} favoreça \textbf{uma} nova concepção educativa.
Refletir sobre os saberes, práticas e vivências através de exercícios e avaliações aplicados ao ensino médio, nos levam a refletir como todas as ciências são sociais, como lemos em Santos ( 2011, p. 89 ), quando a natureza é transformada num artefato global, pois a cultura assim age, passando de “ artefato intrometido num mundo da natureza à expressão da conversão da natureza em artefato total ”.
E a natureza \textbf{enquanto} objeto de conhecimento e entidade cultural possibilitou esta transformação, portanto as ciências ditas naturais são também humanas e sociais.

% matched lemmas: abordagem, atual, autor, basear, compromisso, conseguir, contemporâneo, contextualização, dialético, direção, discente, dizer, docência, educar, enquanto, equidade, erro, exigir, fomentar, forte, fragmentação, fundamento, hegemônico, historicamente, humanidade, interdisciplinaridade, literário, mero, metodológico, método, percepção, permear, poderoso, posicionamento, preciso, que, remeter, revelar, sujeito, século, sólido, tarefa, teoria, teórico, tópico, um, ver, viver, vontade
\end{textsample}
